\chapter{Metodología y Análisis del Sistema Propuesto}

El análisis del estado del arte presentado en el capítulo anterior evidenció las 
limitaciones de las herramientas comerciales y académicas existentes para el diagnóstico 
dermatológico asistido por inteligencia artificial, particularmente en lo referente a la 
falta de adaptación a los fototipos de piel oscura, los altos costos de suscripción, 
la dependencia de conectividad a internet y la ausencia de modelos explicativos (XAI). 
A partir de estas brechas identificadas, el presente capítulo describe la solución 
propuesta para DermaUH versión 3.0. En primer lugar, se define la metodología de desarrollo
de software que guiará el proyecto, justificando la elección de un enfoque ágil adaptado 
a los requisitos del sistema. A continuación, se presenta 
la especificación de requisitos funcionales y no funcionales, estableciendo las 
capacidades que debe ofrecer la plataforma y las condiciones de calidad que debe cumplir. 
Finalmente, se aborda el análisis y diseño del sistema, incluyendo el modelado del negocio 
mediante diagramas de casos de uso UML, el diseño arquitectónico (capas, cliente-servidor, 
MVC) acompañado de su diagrama de componentes, y el diseño detallado con diagramas de 
secuencia, diagrama de clases UML, y el modelo de base de datos (diagrama entidad-relación 
y modelo relacional normalizado). De esta forma, se sientan las bases técnicas para la 
implementación de una solución que responda a las necesidades específicas de los 
dermatólogos cubanos.

\section{Metodolog\'ia}
Una metodología de desarrollo de software es un marco de trabajo que se
usa para estructurar, planificar y controlar el proceso de desarrollo de
sistemas de información. Las metodologías de Desarrollo de Software (DS.) han experimentado un
proceso histórico y evolutivo que inicia en los años 40 con la aparición de las
primeras computadoras. Hoy en d\'ia existen dos tipos de metodolog\'ias, tradicionales y \'agiles. \cite{metodology}

\begin{enumerate}
    \item[1-] \textbf{Metodolog\'ias Tradicionales:} Centran su atenci\'on en llevar una documentaci\'on exhaustiva de todo el proyecto, planificaci\'on y control del mismo. Imponen una disciplina rigurosa, con el fin de lograr un software m\'as eficiente. Algunas de estas metodolog\'ias son:
    \subitem 1.1- \textbf{Cascada:} Consiste en fases secuenciales donde cada una debe completarse antes de pasar a la fase siguiente. Es muy f\'acil de entender y gestionar, es ideal cuando los requisitos son estables y bien conocidos, sin embargo es muy r\'igida antes cambios y el cliente solo ve el producto final.
    \subitem 1.2- \textbf{V-Model:} Es una variante del Cascada que enfatiza la verificaci\'on y validaci\'on (V\&V), cada fase de desarrollo tiene una fase de prueba correspondiente. Presenta una alta calidad y trazabilidad, adem\'as las pruebas se planifican desde el inicio. Por otro lado sigue siendo r\'igido como Cascada y no es apto para requisitos cambiantes.
    \subitem 1.3- \textbf{Modelo Espiral:} Propuesto por Barry Boehm, combina Cascada con gestión de riesgos mediante iteraciones que aumentan en detalle. Este enfoque permite iteraciones controladas y es bueno para proyectos grandes con incertidumbre t\'ecnica, sin embargo es complejo de gestionar y puede ser costoso en cuanto a recursos.

    \item[2-] \textbf{Metodolog\'ias \'Agiles:} Generalmente son procesos \textit{Incrementales} (entregas en ciclos cortos y r\'apidos), \textit{Cooperativo} (clientes y desarrolladores trabajan constantemente con una comunicación muy fina y constante), Sencillo (el método es fácil de aprender y modificar para el equipo) y finalmente Adaptativo (capaz de permitir cambios de último momento). Estas metodologías ponen de relevancia que la capacidad de respuesta a un cambio es más importante que el seguimiento estricto de un plan \cite{metodologiasAgiles}. Entre las m\'as conocidas se pueden apreciar:
    \subitem 2.1- \textbf{Scrum:} Se basa en ciclos de trabajos cortos llamados \textit{Sprints} que tienen una duraci\'on de entre 2 y 4 semanas. Es ideal para proyectos complejos que requieren flexibilidad y entregas r\'apidas. \cite{scrumAbstract}
    \subitem 2.2- \textbf{Kanban:} Kanban se centra en visualizar el flujo de trabajo y limitar el trabajo en progreso (WIP, por sus siglas en inglés). No utiliza iteraciones de tiempo fijo como Scrum, sino que se basa en un flujo continuo de tareas. Es ideal para equipos que buscan mejorar la eficiencia y la gestión de tareas, pero puede ser menos estructurado que Scrum.
    \subitem 2.3- \textbf{Extreme Programming (XP):} XP se enfoca en las prácticas de ingeniería de software para producir código de alta calidad y adaptarse a requisitos cambiantes. Es especialmente adecuado para proyectos con requisitos muy volátiles. Promueve prácticas como la programación en parejas (dos desarrolladores trabajan en el mismo código), el desarrollo guiado por pruebas (TDD), la integración continua y las refactorizaciones frecuentes para mejorar el diseño del código. 
\end{enumerate}

Para el desarrollo de DermaUH 3.0 se ha decidido utilizar la metodología Scrum, debido a que es una metodología ágil que 
se adapta bien a proyectos complejos y con requisitos cambiantes, como es el caso de DermaUH. Además, Scrum permite entregas rápidas y 
frecuentes, lo que facilita la retroalimentación constante del cliente y la adaptación a sus necesidades.
Para ello se dividi\'o el proceso de desarrollo en seis sprints, con una duraci\'on promedio de entre dos y tres semanas cada uno.
Al final de cada sprint se realiz\'o una evaluaci\'on y validaci\'on de los resultados obtenidos. Los sprints en los que fue dividido el proceso
son los siguientes:

\begin{enumerate}
    \item Desarrollo de toda la secci\'on de Registro e Inicio de Sesi\'on de los usuarios. Definir roles y permisos de cada uno.
    \item Implementar un servicio de correo, el cual notifica al usuario administrador de nuevas solicitudes y a los usuarios les notifica la respuesta de su solicitud.
    \item Desarrollar las funcionalidades b\'asicas de la Aplicaci\'on, gesti\'on de pacientes, lesiones y diagn\'osticos.
    \item Entrenar los nuevos modelos a integrar con el nuevo conjunto im\'agenes entregado por las organizaciones de la salud.
    \item Desarrollar el servicio de consulta de Imagen y vincular el modelo a la aplicaci\'on DermaUH 3.0
    \item Realizar pruebas de rendimiento y ejecuci\'on con doctores para posteriormente lanzar la aplicaci\'on.
\end{enumerate}

\section{Levantamiento de Requerimientos}
El levantamiento de requerimientos es el proceso mediante el cual se identifican, 
documentan y acuerdan las necesidades, expectativas y restricciones que debe cumplir un 
sistema de software. Incluye tanto los requerimientos funcionales (qué debe hacer el sistema) 
como los no funcionales (cómo debe hacerlo: rendimiento, seguridad, usabilidad, etc.). 
Este proceso involucra a los stakeholders (clientes, usuarios finales, desarrolladores) y 
constituye la base sobre la que se construye todo el proyecto.

Los requerimientos son importantes porque actúan como el contrato implícito entre el equipo de desarrollo y los interesados. 
Unos requerimientos claros y bien gestionados permiten: definir el alcance realista del proyecto, 
evitar malentendidos y retrabajos costosos, establecer criterios objetivos de validación y aceptación y
priorizar funcionalidades según el valor que aportan. La falta de una buena gestión de requerimientos 
es una de las causas más frecuentes de fracaso en proyectos de software, ya que conduce a entregas 
que no satisfacen las necesidades reales, sobrecostes y demoras.

\subsection{Requerimientos Funcionales}
Seg\'un un art\'iculo de Michel Arias Chaves, para la revista InterSede \cite{requerimientosFuncionales}
los requerimientos funcionales son los que definen las funciones que el sistema ser\'a capaz de realizar,
describen las transformaciones que el sistema realiza sobre las entradas para producir salidas.Es
importante que se describa el ¿Qué? y no el ¿Cómo? se deben hacer esas transformaciones. Estos
requerimientos al tiempo que avanza el proyecto de software se convierten en los algoritmos, la lógica y
gran parte del código del sistema. Siguiendo este concepto, los requerimientos funcionales de DermaUH 3.0 son:
\begin{enumerate}
    \item[RF1] El usuario autenticado puede hacer \textbf{CRUD} sobre las entidades Paciente, Lesi\'on y diagn\'ostico.
    \item[RF2] Los usuarios podr\'an obtener informaci\'on de lesiones registradas, solo de sus propios pacientes, podr\'an acceder a otras lesiones si estas se hacen p\'ublicas o se comparten con el mismo.
    \item[RF3] El usuario puede enviar una imagen dermatosc\'opica o cl\'inica al sistema y recibir una sugerencia de clasificaci\'on automatizada con nivel de confianza, la cual debe ser validada por un especialista m\'edico antes de cualquier decisi\'on cl\'inica.
    \item[RF4] Solamante el administrador tiene permiso para eliminar cuentas y activar usuarios en el sistema y consultar detalles de cualquier cuenta.
    \item[RF5] Cualquier persona puede enviar una solicitud de registro a la aplicaci\'on, la cual deber\'a aportar datos profesionales como el hospital y el carnet de identidad para despu\'es ser\'a aprobada por el Administrador.
\end{enumerate}

\subsection{Requerimientos No Funcionales}
los requerimientos no funcionales tienen que ver con características que de una u otra
forma puedan limitar el sistema, como por ejemplo, el rendimiento (en tiempo y espacio), interfaces de
usuario, fiabilidad (robustez del sistema, disponibilidad de equipo), mantenimiento, seguridad,
portabilidad, estándares, etc. Siguiendo esta l\'inea se concluye que los requerimientos no funcionales
de DermaUH son:
\begin{enumerate}
    \item[RNF1] La autenticaci\'on debe ser mediante email y password generando un JSON Web Token (JWT) el cual definir\'a los permisos del usuario que hace inicio de sesi\'on. El JWT a su vez debe tener un tiempo de expiraci\'on y ser renovable 
    \item[RNF2] La comunicación entre cliente y servidor debe estar cifrada mediante HTTPS.
    \item[RNF3] El tiempo de operaci\'on de una petici\'on y la generaci\'on de la respuesta debe ser inferior a dos segundos con internet estable.
    \item[RNF4] La interfaz web debe ser responsive (adaptable para m\'oviles, tabletas y ordenadores).
    \item[RNF5] Cualquier respuesta a una solicitud, sea positiva o de error debe mostrarse con mensajes claros a los usuarios.
\end{enumerate}

\section{Análisis y Diseño}

Una vez definidos los requisitos funcionales y no funcionales que debe cumplir el sistema, 
el siguiente paso consiste en modelar su estructura y comportamiento desde una perspectiva técnica. 
Este apartado presenta los artefactos de diseño que guiarán la implementación de la versión 3.0 de 
DermaUH. Se abordar\'an temas como: casos de uso que aparecen en la aplicac\'on, arquitectura 
de software seleccionada así como estructura, dise\~no y modelaci\'on de las entidades de 
la base de datos. 

\subsection{Diagrama de Casos de Uso}
Los diagramas de casos de uso se utilizan durante la fase de análisis de un proyecto para
identificar y dividir la funcionalidad del sistema. Normalmente contienen: casos de uso,
actores y relaciones entre ellos: de asociación, de dependencia y/o de generalización\cite{diagramaUML}. 
En otras palabras es una forma de esquematizar y representar los requerimientos funcionales planteados
en el ep\'igrafe anterior.

\begin{itemize}
    \item \textbf{Casos de Uso:} Los casos de uso describen el comportamiento del sistema cuando uno de los actores envía un estímulo concreto.
    \item \textbf{Actores:} Los actores interpretan papeles que pueden ser representados por los usuarios del sistema. Dichos usuarios pueden ser humanos, otros ordenadores o incluso otros sistemas de software 
\end{itemize}

A continuaci\'on se mostrar\'an los diagramas de casos de uso, realizados con la herramienta \href{https://www.plantuml.com/plantuml/uml/}{PlantUML}
debido a que permite mediante su propio lenguaje \textit{plantuml} crear im\'agenes de diagramas
de casos de uso. Normalmente en un diagrama de casos de uso
se visualizan todos los actores y los casos de uso, por simplicidad y para que sea m\'as entendible 
para el lector, se separaron los diagramas por actores.

\newpage

\begin{figure}[ht] 
    \centering
    \includegraphics[width=1.0\textwidth]{figures/diagrama_casos_uso_admin.png}
    \caption{Diagrama UML del actor Admin} 
    \label{adminUML}
\end{figure}

La figura \ref{adminUML} muestra las acciones permitidas cuando el usuario entra al sistema 
con permisos de administrador. A continuaci\'on una breve explicaci\'on de cada una:

En la pestaña \textit{Usuarios}:
\begin{itemize}
    \item \textbf{Listar Cuentas de Usuarios:} Para acceder a eliminar o ver detalles, el administrador primero debe consultar la lista de usuarios
    \item \textbf{Eliminar Cuenta del Usuario:} Esta opci\'on permite eliminar la cuenta de un usuario de la aplicaci\'on, sin importar los permisos del usuario ni el tiempo que lleve activo en el sistema.
    \item \textbf{Ver Detalles del Usuario:} Esta opci\'on le muestra al administrador todos los datos personales que el usuario ingres\'o en el formulario de registro, como su nombre completo, hospital al que pertenece, correo entre otros. En caso del usuario ingresar como residente, muestra tambi\'en el correo del doctor encargado de tutorearlo.
\end{itemize}

En la pestaña \textit{Solicitudes}:
\begin{itemize}
    \item \textbf{Rechazar Solicitud de Usuario:} Con esta opci\'on el administrador rechaza, lo cual es sin\'onimo de eliminar, al menos en este contexto, la solicitud de registro enviada por un usuario.
    \item \textbf{Aceptar Solicitud de Usuario:} Permite al administrador aceptar la solicitud de registro de un usuario y este ya podrá utilizar la aplicación en dependencia de los permisos asignados por el administrador.
    \item \textbf{Asignar Roles al Usuario:} Una vez aceptada la solicitud de registro, el adminsitrador debe asignarle permisos a dicho usuario. Los permisos permitidos son administrador y/o doctor.
    \item \textbf{Enviar correo:} La acción de aceptar o rechazar una solicitud incluye automáticamente el envío de un correo al solicitante, por lo que se modeló como una relación de inclusión (<<include>>)
\end{itemize}

El usuario con permisos de Administrador es el único que tiene poder y potestad
para manejar todo lo relacionados con las cuentas de los usuarios y las solicitudes de registro.
A continuación se mostrará que ocurre en caso del usuario tener permisos de doctor.

\begin{figure}[ht] 
    \centering
    \includegraphics[width=0.8\textwidth]{figures/diagrama_casos_uso_doctor.png}
    \caption{Diagrama UML del actor Doctor} 
    \label{doctorUML}
\end{figure}

\begin{itemize}
    \item \textbf{Registrar Nuevos Pacientes:} En el panel principal de la aplicaci\'on, arriba a la derecha, aparece un bot\'on para registrar pacientes nuevos en el sistema. Al presionarlo se muestra un formulario que debe ser rellenado con informaci\'on personal del paciente, nombre, edad, sexo, profesi\'on etc. Una vez rellenado el formulario, se presiona el bot\'on de finalizar registro y el paciente queda registrado en el sistema. 
    \item \textbf{Ver detalles de un Paciente:} Los pacientes son representados como tarjetas en el panel principal. A la derecha de cada tarjeta se muestra un bot\'on con la etiqueta \textit{Detalles} que permite al doctor visualizar los datos personales del paciente, separada en tres secciones, informaci\'on personal, informaci\'on \'etnica y direcci\'on.
    \item \textbf{Ver lesiones de un Paciente:} Similar a la opci\'on anterior, a la derecha del paciente se muestra el bot\'on \textit{Lesiones}. Este le permite al doctor observar el registro de lesiones del paciente en cuesti\'on.
    \item \textbf{Registrar nuevas lesiones:} Se puede llegar a esta opcio\'on de dos formas diferentes, cuando se registra un nuevo paciente y se presiona el bot\'on finalizar registro, aparece un nuevo formulario para registrar la neuva lesi\'on y la otra es al acceder al listado de lesiones de un paciente se muestra un bot\'on para agregar una lesi\'on nueva. 
    \item \textbf{Eliminar un Paciente:} Similar a ver lesiones de un paciente y ver detalles se accede a esta opci\'on mediante un bot\'on ubicado a la derecha de la tarjeta del paciente.
    \item \textbf{Listar y buscar pacientes mediante filtros:} En la parte superior del panel principal aparece una barra de b\'usqueda, donde el doctor podr\'a introducir una consulta y la aplicaci\'on filtrar\'a a los pacientes por la consulta del doctor. Adem\'as aparece una opci\'on para ordenar los pacientes por alg\'un campo.
    \item \textbf{Compartir lesiones con otros doctores:} Con esta opci\'on, el doctor tiene la posibilidad de compartir una lesi\'on con otros colegas del personal de la salud. Los doctores escogidos para compartir tendr\'an la posibilidad de ver solamente el nombre y apellido del paciente y las im\'agenes de la lesi\'on.
    \item \textbf{Solicitar diagn\'ostico a trav\'es de la imagen:} Esta opci\'on permite al doctor tomar una fotograf\'ia, ya sea con el m\'ovil o con el dermatoscopio y solicitar un diagn\'ostico aproximado de confianza proporcionado por un modelo de ML especializado en procesamiento de im\'agenes.
\end{itemize}

\subsection{An\'alisis y propuesta arquitect\'onica}
Una vez modelados los casos de uso que describen el comportamiento funcional del 
sistema desde la perspectiva de los actores, el siguiente paso consiste en definir 
la estructura organizativa que soportará dichas funcionalidades. La arquitectura de 
software se refiere a la disposición conceptual y estructural que subyace en el diseño y 
desarrollo de sistemas, estableciendo los componentes, sus relaciones y los principios 
que guían su evolución. Su importancia radica en proporcionar un marco de trabajo para 
cumplir con los requisitos funcionales y no funcionales (rendimiento, seguridad, 
escalabilidad), así como para tomar decisiones tecnológicas fundamentadas\cite{architecture}. En este 
apartado se realizar\'a un an\'alisis de las arquitecturas de software m\'as conocidas
atendiendo a su capacidad de escalabilidad, implementaci\'on para solucionar el problema en cuesti\'on
y los conocimientos adquiridos en la carrera de Licenciatura en Ciencias de la Computaci\'on.

\subsubsection{Arquitectura Hexagonal (Puertos y Adaptadores)}
La arquitectura hexagonal, también conocida como arquitectura de puertos y adaptadores, 
fue propuesta por Alistair Cockburn\cite{hexagonal1}. Su principio fundamental es aislar el dominio 
(lógica de negocio) de los detalles técnicos externos como interfaces de usuario, bases 
de datos o servicios de mensajería. El dominio se ubica en el centro de un hexágono 
imaginario; alrededor de él se definen \textbf{puertos} (interfaces que expresan las 
capacidades del dominio) y \textbf{adaptadores} (implementaciones concretas que conectan 
cada puerto con una tecnología externa). 

Por ejemplo, un puerto podría ser \texttt{RepositorioPaciente} (interfaz para guardar y 
recuperar pacientes). Un adaptador concreto sería \texttt{RepositorioPacientePostgreSQL}, 
que implementa esa interfaz usando SQL. Gracias a esta separación, el dominio nunca depende
de la base de datos ni del framework web, sino que las dependencias apuntan hacia adentro.
Esto facilita probar el dominio de forma aislada (con adaptadores de prueba), 
cambiar de base de datos sin tocar la lógica de negocio, y retrasar decisiones 
tecnológicas.

La arquitectura hexagonal es especialmente útil en sistemas con lógica de negocio compleja
y larga vida útil, donde se esperan cambios en los detalles tecnológicos. No obstante, 
puede resultar excesiva para prototipos o aplicaciones muy simples.

\begin{figure}[ht] 
    \centering
    \includegraphics[width=0.5\textwidth]{figures/hex-ports-adapters-en.png}
    \caption{Esquema Arquitectura Hexagonal} 
    \label{hexagonalschema}
\end{figure}

\subsubsection*{Onion Architecture}
La Onion Architecture (o en capas de cebolla) fue implementada y presentada por Jeffrey
Palermo en 2008. Como todo patr\'on de arquitectura, este hace hincapi\'e en la divisi\'on
de funciones y el desarrollo de componentes interconectados. Desde que vio la luz en el a\~no
2008 ha ganado popularidad entre los ingenieros de software que buscan crear aplicaciones
escalables, mantenibles y probadas. 

Seg\'un el art\'iculo Onion Architecture Used in Software Development\cite{onion} publicado
en \href{https://www.researchgate.net}{ResearchGate}, plataforma dedicada a la publicaci\'on 
de art\'iculos de ind\'ole cient\'ifico: En la Onion Architecture, Los componentes del 
software se organizan en capas concéntricas mediante la arquitectura de cebolla, 
un patrón de arquitectura en capas en el que cada capa depende exclusivamente de la capa 
inmediatamente posterior. La lógica central de la aplicación se ubica en el centro de la 
arquitectura, rodeada por varias capas que la encierran y la protegen del exterior, razón 
por la cual este patrón se denomina ``cebolla''. La principal idea es que las aplicaciones
sean divididas en cuatro capas:

\begin{enumerate}
    \item \textbf{Dominio:} Esta capa actúa como el cerebro de la aplicación y representa los conceptos, entidades y reglas de negocio. No puede depender de otras capas, ya que es el centro de la arquitectura y tecnol\'ogicamente independiente.
    \item \textbf{Infraestructura:} En esta capa se implementan las interfaces definidas por la capa de dominio. Adem\'as proporciona los elementos de infracestructura como mensajer\'ia, registro de eventos y acceso a datos.
    \item \textbf{Aplicaci\'on:} La lógica de negocio relevante para la aplicación se encuentra en la capa de aplicación, que también actúa como enlace entre las capas de dominio y presentación. Convierte los casos de uso y las reglas de negocio de la capa de dominio en acciones que la capa de presentación puede comprender.
    \item \textbf{Presentaci\'on:} Esta capa ofrece la interfaz para que el usuario interactúe con la aplicación. Una aplicación web, una aplicación de escritorio o una aplicación móvil pueden conformar esta capa.
\end{enumerate}

\begin{figure}[ht] 
    \centering
    \includegraphics[width=0.5\textwidth]{figures/onion.png}
    \caption{Esquema Onion Architecture} 
    \label{onionArchitecture}
\end{figure}
\newpage

La figura~\ref{onionArchitecture} muestra un esquema de estructura de Onion Architecture, para 
comprender su dise\~no y el concepto de capas externas que ``protegen'' a la capa interna mencionado
anteriormente. Esta arquitectura posee un gran n\'umero de ventajas y desventajas para el desarrollo
de software, sin embargo para evitar la extensi\'on del ep\'igrafe se enunciar\'an unas pocas 
y de manera resumida.

\paragraph{Ventajas:}
\begin{enumerate}
    \item Separaci\'on de responsabilidades: La arquitectura de capas fomenta una clara división de responsabilidades entre las múltiples capas de la aplicación. Esto aumenta la modularidad de la aplicación y facilita su comprensión y mantenimiento.
    \item Testibilidad: Debido a que las capas de la arquitectura de cebolla están conectadas de forma flexible y pueden probarse por separado, escribir pruebas automatizadas para la aplicación resulta más sencillo.
    \item Escalabilidad:Gracias a que las capas de la arquitectura cebolla pueden escalarse independientemente unas de otras, las aplicaciones pueden escalarse horizontalmente. A medida que la aplicación se expande, esto ofrece mayor rendimiento y fiabilidad.
    \item Mantenibilidad: Los cambios en una capa no deberían afectar a las demás, ya que cada capa se encarga de una tarea específica. Esto simplifica la actualización y el mantenimiento de la aplicación a lo largo del tiempo.
\end{enumerate}

\paragraph{Desventajas:}
\begin{enumerate}
    \item Mayor complejidad: Incluso en proyectos pequeños, la arquitectura cebolla puede dificultar el desarrollo de una aplicación. La aplicación puede volverse más difícil de comprender y modificar debido a los niveles e interfaces adicionales.
    \item Curva de aprendizaje: La comprensión y aplicación adecuadas de la arquitectura pueden requerir más esfuerzo y formación para quienes no estén familiarizados con la Arquitectura Cebolla.
    \item Sobrediseño: En proyectos pequeños, la arquitectura de cebolla puede, en ocasiones, estar sobrediseñada, lo que añade complejidad y sobrecarga.
    \item Sobrecarga de rendimiento: La sobrecarga de rendimiento puede deberse a las capas e interfaces adicionales de la arquitectura de cebolla, especialmente en aplicaciones que requieren alto rendimiento o procesamiento en tiempo real.
\end{enumerate}

\subsubsection*{Clean Architecture}
La Arquitectura Limpia (Clean Architecture) es un enfoque de dise\~no de software propuesto
por Robert C. Martin que busca la independencia de los sistemas con respecto a frameworks,
bases de datos y detalles de implementaci\'on\cite{cleanArchitecture}. A través de la separación 
de responsabilidades, Clean Architecture permite que los sistemas sean adaptables a 
cambios tecnológicos sin alterar la lógica de negocio principal. 

Clean Architecture se basa en conceptos previos como la Arquitectura en Capas, 
la Arquitectura Hexagonal (Ports and Adapters) de Alistair Cockburn y la Arquitectura en 
Cebolla (Onion Architecture) de Jeffrey Palermo, pero con un enfoque más claro en la 
independencia de los componentes y la testabilidad del código. Clean Architecture
se gu\'ia por cuatro principios fundamentales\cite{cleanArchitecturePrincipios} los cuales
ser\'an explicados a continuaci\'on:

\begin{enumerate}
    \item \textbf{Independencia de cualquier framework:} La arquitectura limpia debe poder aplicarse a cualquier sistema, independientemente del lenguaje de programación o las bibliotecas/frameworks que utilice. Las capas deben estar tan bien separadas que puedan sobrevivir individualmente, sin necesidad de elementos externos.
    \item \textbf{Testeable:} Cada módulo, tanto UI, base de datos, conexión API Rest, etc., debe poder probarse individualmente.
    \item \textbf{Interfaz de usuario (UI) independiente:} La interfaz de usuario debe poder cambiar sin interrumpir todo el sistema. Esta capa debería vivir de forma tan independiente como para ser desmontada y sustituida por otra.
    \item \textbf{Base de datos independiente:} Al igual que en el punto anterior, esta capa debe ser tan modular como para agregar múltiples fuentes de datos e incluso múltiples fuentes del mismo tipo de datos.
\end{enumerate}

Como se mencion\'o anteriormente, Clean Architecture toma fundamentos de la arquitectura
Hexagonal y la Onion, ambas explicadas en este ep\'igrafe. Por tanto tiene una divisi\'on
por capas. Existen bibliograf\'ias en las que el esquema para presentar a esta arquitectura
es similar al esquema Onion, (figura~\ref{onionArchitecture}) y en otras se representa mediante
un cono, como se muestra en la figura~\ref{cleanArchitecture}

\begin{figure}[ht] 
    \centering
    \includegraphics[width=0.5\textwidth]{figures/clean.png}
    \caption{Esquema Clean Architecture en forma de Cono} 
    \label{cleanArchitecture}
\end{figure}
\newpage

\begin{enumerate}
    \item \textbf{Entidades:} Las entidades encapsulan las reglas comerciales de toda la empresa. Una entidad puede ser un objeto con métodos o puede ser un conjunto de funciones y estructuras de datos. No importa, siempre y cuando las entidades puedan ser utilizadas por muchas aplicaciones diferentes en la empresa.
    \item \textbf{Casos de Uso:} El software de esta capa contiene reglas comerciales específicas de la aplicación. Encapsula e implementa todos los casos de uso del sistema. Estos casos de uso organizan el flujo de datos hacia y desde las entidades, y dirigen a esas entidades para que utilicen sus reglas de negocios en toda la empresa para lograr los objetivos del caso de uso.
    \item \textbf{Adaptadores:} El software en esta capa es un conjunto de adaptadores que convierten datos del formato más conveniente para los casos de uso y entidades, al formato más conveniente para alguna agencia externa como la Base de Datos o la Web. De manera similar, los datos se convierten, en esta capa, de la forma más conveniente para las entidades y los casos de uso, a la forma más conveniente para cualquier marco de persistencia que se esté utilizando. Es decir, la base de datos. Ningún código dentro de este círculo debería saber nada sobre la base de datos. También en esta capa hay cualquier otro adaptador necesario para convertir datos de algún formato externo, como un servicio externo, al formato interno utilizado por los casos de uso y las entidades.
    \item \textbf{Frameworks y controladores} La capa más externa generalmente se compone de frameworks y herramientas como la base de datos, el framework web, etc. Generalmente, no escribe mucho código en esta capa, aparte del código que se comunica con el siguiente círculo hacia adentro.
\end{enumerate}

\subsection{Modelado de la Base de Datos}

Definida la arquitectura del sistema, el siguiente paso consiste en estructurar los datos que 
persistirán y las relaciones entre ellos. El presente apartado describe el modelo de datos de 
DermaUH 3.0, incluyendo las entidades principales (usuario, paciente, lesión, imagen, 
diagnóstico automático, entre otras), sus atributos, las claves primarias y foráneas, y 
las asociaciones que garantizan la integridad y coherencia de la información. Se presenta 
el Modelo Entidad-Relación, tambi\'en conocido como Diagrama de Entidad Relaci\'on (DER).

\subsubsection{Modelo de Entidad Relación (DER)}

Seg\'un un art\'iculo obtenido del sitio web UNIR, herramienta elaborada, desarrollada
y dise\~nada con el prop\'osito de ayudar a los estudiantes de universidades europeas,
El modelo entidad relación (ERD o modelos ER) es una herramienta que permite representar 
de manera simplificada cómo personas, objetos o conceptos se relacionan entre sí. 
Se utiliza para exponer cómo se organiza la información en una base de datos. 
La creación del modelo entidad relación para su aplicación en el diseño de bases de 
datos se le atribuye a Peter Chen, profesor del MIT, quien publicó en 1976 el documento  
“Modelo entidad-relación: hacia una visión unificada de los datos”. A continuaci\'on se
expone el DER utilizado para la modelaci\'on de la base de datos.


TODO: HACER EL DIAGRAMA CUANDO LA HERRAMIENTA LE DE LA GANA

\subsubsection{Paciente}
El modelo de datos, \textbf{Paciente}, es la entidad m\'as importante, ya que es el n\'ucleo
de una de las principales funcionalidades de DermaUH 3.0. A continuaci\'on una descripci\'on 
de los atributos de dicha entidad.
\begin{itemize}
    \item idPaciente: Llave primaria (PK), identificador \'unico de tipo entero. Es generado auto\'aticamente por el sistema de manera incremental.
    \item name, lastName1, lastName2: Campos para obtener el nombre completo del paciente. Todos son de tipo varchar, con una longitud m\'axima de cien caracteres por cada campo.
    \item yeras, sex, skinColor: Campos relacionados con la informaci\'on \'etnica del paciente. Estos campos su principal funcionalidad es brindar m\'as informaci\'on para obtener un mejor diagn\'ostico asistido por parte de los modelos ML integrados en la nueva versi\'on.
    \item addres, city, state, jobProfession: Campos para reflejar la direcci\'on y profesi\'on laboral del paciente. Todos son de tipo varchar sin l\'imites. Estas propiedades son importantes para saber el nivel exposici\'on de la piel al sol, debido la heterogeneidad del clima en diferentes sitios del pa\'is.
    \item registerData: Campo de tipo Date para tener control de cuando fue que el paciente fue registrado por primera vez en el sistema de DermaUH.
\end{itemize}

\subsubsection{Dermat\'ologo}
Esta es una entidad fuerte, ya que no depende de otras entidades para instanciarse. El modelo
\textbf{Dermat\'ologo}, es el responsable de representar la informaci\'on de los usaurios del 
sistema. En pocas palabras, un dermat\'ologo es un usuario de DermaUH 3.0.
\begin{itemize}
    \item idDermatologo: Llave primaria (PK), identificador \'unico de tipo entero. Es generado auto\'aticamente por el sistema de manera incremental.
    \item name, lastName1, lastName2: Campos para obtener el nombre completo del dermat\'ologo. Todos son de tipo varchar, con una longitud m\'axima de cien caracteres por cada campo.
    \item email: Campo de tipo text, se usa para guardar el email del dermat\'ologo. En el sistema este campo es utilizado como username (nombre de usuario) y como el identificador para hacer inicio de sesi\'on.
    \item ci: Carnet de Identidad del Usuario, campo de tipo varchar.
    \item phoneNumber: Campo para guardar el n\'umero telef\'onico del usuario.
    \item hospital, state, city: Campos de tipo de varchar para guardar el hospital, municipio y provincia respectivamente del usuario.
    \item doctorId: Identificador que se genera seg\'un el hospital, municipio y provincia del dermat\'ologo.
    \item userType: Campo de tipo varchar. Solo hay tres posibles valores para este campo: Admin, Doctor, Resident, es utilizado para manejar los permisos y roles de los usuarios de la aplicaci\'on.
    \item tutorEmail: Campo de tipo varchar, almacena el correo del tutor en caso del usuario ser residente, o estudiante de intercambio. En caso contrario este campo se puede mantener null.
\end{itemize}

\subsubsection{Lesi\'on}
Entidad d\'ebil, modela el comportamiento de las lesiones que puede presentar un paciente. Un
paciente puede presentar muchas una o muchas lesiones, por tanto el modelo \textbf{Lesi\'on}
depende de que exista un paciente. Se configur\'o con eliminaci\'on \textbf{Cascada}, por 
ende si un paciente es eliminado del sistema, se eliminar\'an todas sus lesiones asociadas.
\begin{itemize}
    \item date: Fecha en la que se registra la lesi\'on. Es un campo de tipo date, no editable.
    \item phototype: Campo de tipo entero, solo puede tomar valores en el intervalo $[1;6]$ que se usa para almacenar el fototipo de la piel donde se encuentra ubicada la lesi\'on. El fototipo no es m\'as que la reacci\'on de la piel ante la exposici\'on a la luz solar.
    \item localizarion: Campo de tipo entero, toma valores en el intervalo $[1;10]$ y almacena la parte del cuerpo donde se encuentra la lesi\'on. 
    \item histopathologicalDiagnosis: Campo de tipo bool. Indica si la lesi\'on fue diagnosticada histopat\'ologicamente, en otras palabras, afirma que se utilizaron equipos m\'edicos para obtener el diagn\'ostico.
    \item injuryType: Llave for\'anea (FK) que referencia a la entidad \textbf{Tipo de Lesi\'on}. Es un campo de tipo entero.
    \item patientId: Llave for\'anea (FK) que referencia a la entidad \textbf{Paciente}. Es un campo de tipo entero.
\end{itemize}